% Generated draft: requires mathematical and source review.
\documentclass[11pt]{article}
\usepackage[margin=25mm]{geometry}
\usepackage{fontspec}
\setmainfont{Latin Modern Roman}
\usepackage{amsmath,amssymb,mathtools,unicode-math}
\setmathfont{Latin Modern Math}
\usepackage{enumitem,needspace,xurl,hyperref}
\setlength{\parindent}{0pt}
\setlength{\parskip}{5pt}
\setlength{\emergencystretch}{3em}
\newcommand{\N}{\mathbb{N}}
\newcommand{\Z}{\mathbb{Z}}
\newcommand{\Q}{\mathbb{Q}}
\newcommand{\R}{\mathbb{R}}
\newcommand{\C}{\mathbb{C}}
\newcommand{\F}{\mathbb{F}}
\newcommand{\A}{\mathbb{A}}
\newcommand{\PP}{\mathbb{P}}
\newcommand{\E}{\mathbb{E}}
\newcommand{\eps}{\varepsilon}
\newcommand{\dd}{\,\mathrm{d}}
\begin{document}
\section*{UnsolvedMath ID: 801109}
\textbf{Source problem code:} catalogue entry
\subsection{Definitions and mathematical statement}
A type $\mathrm{II}_1$ factor is an infinite-dimensional finite von Neumann algebra with center $\C 1$ and normalized faithful normal trace. We interpret separability as separability of the predual. A type $\mathrm{II}_1$ factor $M$ is prime if there are no diffuse type $\mathrm{II}_1$ factors $P,Q$ such that
\[
M\cong P\,\overline\otimes\,Q,
\]
where $\overline\otimes$ denotes the spatial von Neumann algebra tensor product.

We interpret an \emph{existential copy} of $N$ in $M$ as a unital normal trace-preserving $*$-embedding $\iota:N\hookrightarrow M$ that is existential in the continuous first-order language of tracial von Neumann algebras. Explicitly, for every quantifier-free formula $\varphi(\bar x,\bar y)$ and every finite parameter tuple $\bar a$ from $N$,
\[
\inf_{\bar b\in N}\varphi^N(\bar b,\bar a)
=
\inf_{\bar c\in M}\varphi^M(\bar c,\iota(\bar a)),
\]
with all variables ranging over the operator-norm ball sorts prescribed by the formula. The supplied record does not specify this language or its convention for existential embeddings; this is the standard model-theoretic interpretation used in this draft.

\textbf{Question.} For every separable type $\mathrm{II}_1$ factor $M$ and every nonamenable biexact type $\mathrm{II}_1$ factor $N$, does the existence of such an existential embedding $\iota:N\hookrightarrow M$ imply that $M$ is prime?

Here nonamenability is understood in the von Neumann algebra sense. The hypothesis that $N$ is \emph{biexact} is retained exactly as supplied. The record gives no definition or convention for biexactness of a von Neumann factor, so this hypothesis requires clarification from the source before the formulation can be considered fully specified.
\subsection{Short English statement}
Must every separable type $\mathrm{II}_1$ factor containing an existential copy of a nonamenable biexact type $\mathrm{II}_1$ factor be prime? This draft uses the usual tracial model-theoretic interpretation of existential embedding; the intended definition of biexactness remains unspecified in the supplied record.
\subsection{Sources}
\begin{itemize}
\item Supplied local UnsolvedMath registry record, permanent ID $801109$, exact source problem code \texttt{catalogue entry}; title: \emph{Ultraproduct $\mathrm{II}_1$ Factors --- Problem 10}. The registry attributes the problem to Elayavalli, with proposed year 2025, and labels it open. This label is supplied source metadata, not an independently verified mathematical status.
\item Elayavalli, \emph{50 Open Problems: Ultraproduct $\mathrm{II}_1$ Factors} (2025), Problem 10, as identified in the supplied provenance: \url{https://arxiv.org/abs/2511.20377}. The record reports access on 2026-10-16 and extraction by \texttt{arxiv\_latex}; these are supplied provenance claims, not a claim of source consultation for this draft.
\item Supplied registry revision: \texttt{d371ee917ac9647fbb06cd6fc529deec6f3dda82}. The record was created and updated on 2026-10-16.
\end{itemize}
\end{document}
